\cotefile{134-2}
\scanpage{106}\mspage{\handwritten{$7'$}}

understood that the``amalgamated sums'' we are considering are of the most common type, -- when $X$ is the finite union of closed subsets $X_i$, with given isomorphisms
\[
X_i \handwritten{\simeq} D_{n(i)}\quad ,
\]
the intersection of any two of these $X_i \handwritten{\unsure{\cap|\wedge}} X_j$ being a union of closed cell\unsure{s} both in $D_{n(i)}$ and in $D_{n(j)}$. (This implies in fact that it is either a closed cell\deleted{\unsure{s}} in both, or the union of two closed cells of same dimensio\unsure{n} $m$ and hence isomorphic to $S_m$, -- a case which will be ruled out anyhow by the triviality condition which follows.) The triviality or ``standartness'' condition is now expressed by demanding that the set of indices $I$ can be totally ordered\inserted{, i.e.``numbered''}, in such a way that we get $X$ by succesively ``attaching'' cells $D_{n(i)}$ \inserted{\handwritten{$(i\unsure{\geq|>}2)$}} to the space already constructed\deleted{\unsure{,}} $X(i-1)$, by a map from a \inserted{sub-}cell \deleted{\unsure{$A_i$ and}}\inserted{\handwritten{of $D_{n(i)}$}}, $S_j^{\handwritten{\varepsilon}}\handwritten{\hookrightarrow}X(i-1)$ $(\handwritten{\varepsilon\in\{}\mathord{\pm}1\handwritten{\}})$, this map of course inducing an isomorphism, more precisely \uline{the} standart isomorphism\inserted{\handwritten{${\unsure{\rho|p|f}}_{n(i'),j}^{\unsure{\varepsilon'|\varepsilon}}$}} of $S_j^{\handwritten{\varepsilon}}\handwritten{\simeq} D_j$ with one\inserted{\handwritten{$S^{\varepsilon'}$}} of the two corresponding cells $S_j^{\handwritten{+}}$, $S_j^{\handwritten{-}}$ in some $X_{i'}\handwritten{\simeq}D_{n(i')}$ . The \deleted{\unsure{trans}} dual translation of this in terms of fiber products in \deleted{\unsure{Ens}} a category $C$ endowed with objects $F_i$ $(i\handwritten{\in}\underline{N})$ and maps $s^i_{\unsure{1|l}},t^i_1$ between these, is clear: for a given set of indices $I$ and map $i\handwritten{\mapsto}n(i):I\handwritten{\to}\underline{N}$, we consider a sub\deleted{set}\inserted{object} of \inserted{$P=$}$\handwritten{\prod_{i\in I}}F_{n(i)}$, which can be described by equality relations between iterated sources and targets \deleted{\unsure{between}} of various components of $u=(u_i)_{i\handwritten{\in}I}$ in $P$, the structure of the set of relations being such that $I$ can be numbered, from $1$ to $N$ say, in such a way that we get in succession $N$\inserted{$-1$} relations on the $N$ components $u_i$ \inserted{$(2\handwritten{\leq}i\handwritten{\leq}N)$} re\deleted{la}\handwritten{\unsure{spec}}tively, every relation being of the type $f(u_i)=g(u_{i'})$, with $f$ and $g$ being iterated source or target-maps, and $i'\handwritten{<}i$. \handwritten{(Whe\deleted{\unsure{ther}}ther source or target depending in obvious way on the two signs $\varepsilon,\varepsilon'$.)}

% archive sheet number: 105
\scanpage{107}

\hfill 21.2.\\

\handwritten{8.} Returning to theamalgamated sum $X=\handwritten{\bigcup_i}X_i$, the cellular decompositions of the components $X_i\handwritten{\simeq}D_{n(i)}$ define a cellular decomposition of $X$, whose set of cells with incidence relation forms a finite ordered set $K$, finite union of a family of subsets $(K_i)_{i\handwritten{\in}I}$, with given isomorphisms\deleted{\unsure{ of}}
\[
f_i:K_i\handwritten{\simeq}\underline{J}_{n(i)}\qquad(i\handwritten{\in}I)\quad ,
\]
where for every index $n\handwritten{\in}\underline{N}$, $\underline{J}_n$ denotes the ordered set of the $2n+1$ cells $S_j^{\handwritten{\varepsilon}}$ $(0\handwritten{\leq}j\handwritten{\leq}n-1,\handwritten{\varepsilon\in\{}\pm1\handwritten{\}}),D_n$ of \deleted{\unsure{$D_n$}}the pertinent cellular decomposition of $D_n$. We may without loss of generality assume there is no\deleted{\unsure{n}} inclusion relation between the $K_i$, moreover the standartness condition described above readily translates into a condition on this structure $K$, $(K_i)_{i\handwritten{\in}I}$, $(f_i)_{i\handwritten{\in}I}$, and implies that for $i,i'\handwritten{\in}I$, $K_i\handwritten{\to}K_{i'}$ is a``closed'' subset in the two ordered sets $K_i,K_{i'}$ (namely contains with any element $x$ the elements smaller than $x$), and moreover isomorphic (for the induced order) to some $\underline{J}_n$. Thus the category $B_{\unsure{o|0}}$ \deleted{is} can be viewed as the category of such ``standart ordered sets'' (with the extra structure on these just said), and the category $C_o$ can be defined most simply as the dual category $B_o^{o}$. (NB the definition of morphisms in $B_o$ is clear I guess \ldots) I believe the category $B_o$ is stable under amalgamated sums $X\handwritten{\amalg}_ZY$, provided \inserted{however} we insist that the empty structure $X$ is \uline{not} allowed -- otherwise we have to \deleted{insist} restrict to amalgamated sums with $Z\handwritten{\ne\varnothing}$. It seems finally more convenient to exclude the empty structure in $B_o$, i.e. to exclude final element from $\underline{C}_o$, for the benefit of being able to state that $\underline{C}_o$ (and \uline{all categories} $\underline{C}_n$\handwritten{*}) \deleted{should be} are stable under binary amalgamated sums, and that the functors $\underline{C}_n\handwritten{\to}C$ (and ultimately $\underline{C}_{\handwritten{\infty}}\handwritten{\to}C$) we are interested in are just those commuting to arbitrary binary amalgamated sums (without \deleted{\unsure{refere}} awkwa\handwritten{\unsure{rd}} reference to the objects $\underline{F}_i$ and the iterated source and target maps between them \ldots).
% slanted
\margin{\handwritten{This is obviously false, cf \unsure{PS|P.S} p.12}}

\handwritten{9.} The category $\underline{C}_o$ being fairly well understood, remains to complete the construction by the inductive step, passing from $\underline{C}_n$ to $\underline{C}_{n+1}$. The main properties I have in mind therefore, for the sequence of categories $\underline{C}_n$ and their limit $\underline{C}_{\handwritten{\infty}}$, are the following two.

A) For any $\unsure{K}\handwritten{\in}\operatorname{Ob}(\underline{C}_{\handwritten{\infty}})$ $(\unsure{=}\operatorname{Ob}(\underline{C}_o))$, and any two arrows in $\underline{C}_{\handwritten{\infty}}$
\[
 f,g: \quad K\handwritten{\rightrightarrows}F_i\quad ,
\]
with $i\handwritten{\in}\underline{N}$\deleted{\unsure{,}}, and such that either $i=0$, or the equalities
\[
 s_{\unsure{1|l}}^i f=s_1^i g \quad , \quad t_1^i f=t_1^i g \tag*{\handwritten{(1)}}
\]
hold (case $i\handwritten{\geq}1$), there exists $h:K\quad F_{i+1}$ such that
\[
 s_1^{i+1}h=f\quad ,\quad t_1^{i+1}h=g\quad . \tag*{\handwritten{(2)}}
\]

B) For any $n\handwritten{\in}\underline{N}$, the category $\underline{C}_{n+1}$ is deduced from $\underline{C}_n$ by keeping the same objects, and just adding new arrows \deleted{\unsure{as}} $h$ as in A), with $f,g$ arrows in $\unsure{C|\,\underline{C}_n}$.

% archive sheet number: 106
\scanpage{108}\mspage{\handwritten{$8'$}}

The expression ``deduces from'' in B) means that we are adding arrows $h:K\handwritten{\to}F_i$ (each with preassigned source and target in $\underline{C}_n$), with as ``new axioms'' on \inserted{the bunch of} these uniquely the two relations $(\deleted{\unsure{a}}2)$ of \deleted{\unsure{B}}A), the category $\underline{C}_{n+1}$ being deduced from $\underline{C}_n$ in an obvious way, as the solution of a universal problem with\handwritten{in} the category of all categories where binary amalgamated products exist, -- and ``maps'' between these being functors which commut\unsure{e} to those fibered products.
% slanted
\margin{\handwritten{\unsure{unaccurate|inaccurate} -- amalgamated sums exist in $\underline{C}_{\unsure{o|0}}$ only, and \uline{these} should be \unsure{respected}\ldots}}
In practical terms, the arrows of $\underline{C}_{n+1}$ are those deduced from the arrows in $\underline{C}_n$ and the ``new'' arrows $h$, by combining formal operations of composing arrows by $v\handwritten{\circ}u$, and taking \inserted{\handwritten{(binary)}} amalgamated products of arrows. \deleted{\unsure{Of}}

NB Of course the condition $(1)$ in A) is necessary for the existence of an $h$ satisfying $(2)$. That it is sufficient too can be viewed as an extremely strong,\deleted{\unsure{version}} ``universal'' version of coherence conditions, concerning the various structures introduced on a $\handwritten{\infty}$-groupoid. Intuitively, it means that whenever we have \deleted{a}\inserted{two} way\handwritten{s} \deleted{to}of associat\unsure{ing} to a \inserted{finite} family \deleted{of} $(u_i)_{i\handwritten{\in}I}$ of objects of an $\handwritten{\infty}$-groupoid, $u_i\handwritten{\in}F_{n(i)}$, \deleted{\unsure{and such}} subjected to a ``standart'' set of relations on the $u_i$'s, an element of some $F_n$, in terms of the $\handwritten{\infty}$-groupoid structure only, then we have automatically \deleted{an isomorphism}\inserted{a ``homotopy''} between these built in in the very structure of the $\handwritten{\infty}$-groupoid, provided it makes at all a sense to ask for one (namely provided condition $(1)$ holds if $n\handwritten{\geq}1$). I have the feeling moreover that conditions (A) and (\unsure{B|3}) (plus the relation $\underline{C}_{\handwritten{\infty}}=\handwritten{\varinjlim}\underline{C}_n$) is all what will be ever needed, when using the definition of the structure species, -- plus of course the description of $\underline{C}_o$, and the \inserted{implicit} fact that the categories $\underline{C}_n$ are stable under binary fiber products and \inserted{\handwritten{that}} the inclusion functors commute to these.
% slanted
\margin{\handwritten{unaccurate -- see above!}}
Of course, the category which really interests us is $\underline{C}_{\handwritten{\infty}}$, the description of the intermediate $\underline{C}_n$\handwritten{'s} is merely technical \deleted{\unsure{but}} -- the main point is that there should exist an increasing sequence $(C_n)$ of subcategories of $\underline{C}_{\handwritten{\infty}}$, having same objects\handwritten{[}and ``same'' fiberproducts\handwritten{]}, such that $\underline{C}_{\handwritten{\infty}}$ should be the limit (i.e. every arrow in $C_{\handwritten{\infty}}$ should belong \deleted{\unsure{t}} to some $\underline{C}_n$), and such that the passage from $\underline{C}_n$ to $\underline{C}_{n+1}$ should satisfy (B). It is fairly obvious that these conditions alone do by no means caracterise $\underline{C}_{\handwritten{\infty}}$ up to equivalence, and still less the sequence of its subcategories $\underline{C}_n$. The point I wish to make though, before pursuing with a proposal of an explicit description, is that \uline{this ambiguity is in the nature of things}. \handwritten{R}oughly saying, two different mathematicians, working independently on the conceptual problem I had in mind, \deleted{\unsure{aski}}\unsure{assu}ming they both wind up with some explicit definition, will almost certainly get non equivalent definitions -- namely with non equivalent categories of (set-valued, say) $\handwritten{\infty}$-groupoids ! And, secondly and as importantly,

% archive sheet number: 107
\scanpage{109}\mspage{\handwritten{\unsure{9}}}

that \uline{this ambiguity however} is an \uline{irrelevant one.} To make this point a little clearer, I could say that a third mathematician, informed of the work of both, will readily think\unsure{.|-|\,}out a functor or rather a pair of functors, associating to any structure of Mr. X one of Mr. Y and conver\deleted{\unsure{s}}sely, in such a way that by composition of the two, we will associate to a X-structure \handwritten{(}$T$ say\handwritten{)}another $T'$, which \deleted{of course} will not be isomorphic to $T$ of course, but endowed with a canonical \deleted{\unsure{isom}} $\handwritten{\infty}$-equivalence (in the sense of Mr. X) $T\handwritten{\underset{\infty}{\simeq}}T'$, and the same on the Mr. Y side. Most probably, a fourth mathematician, faced with the same situation as the third, will get \deleted{a different}\inserted{his own} pair of functors to reconciliate Mr X and Mr Y, \deleted{but}\inserted{which} very probably \deleted{his pair of functors} wont be equivalent \deleted{to} \handwritten{(}I mean isomorphic\handwritten{)} to the previous one. Here however a fifth mathematician, informed about this new perplexity, will probably show that the two Y-structures $U$ and $U'$, associated by his two colleagues to \deleted{the} an X-structure $T$, while not isomorphic alas, admit however a canonical $\handwritten{\infty}$-equivalence between $U$ and $U'$\inserted{ (in the sense of the Y-theory)}. I could go on with a sixth mathematici\deleted{a}an, confronted with the same perplexity as the previous one, who winds up with another \inserted{\handwritten{$\infty$}-}equivalence between $U$ and $U'$ (without being informed of the work of the fifth), and a seventh reconciling them by discovering an $\handwritten{\infty}$-equivalence between these equivalences. The story of course is infinite, I better stop with seven mathematicians, a fair number nowadays to allow themselves getting involved with foundational matters \ldots\ There should be a mathematical statement though resuming in finite terms this infinite story, but in order to write it down I guess a \unsure{m}inimum amount of conceptual work, in the context of a given notion of $\handwritten{\infty}$-groupoids satisfying the desiderata (A) and (B), should be done, and I am by no means sure I will go through this, not in this letter anyhow.

\handwritten{10.} Now in the long last the explicit description I promised of $\underline{C}_{n+1}$ in terms of $\underline{C}_n$. As a matter of fact, I have a handful to propose ! O\deleted{\unsure{n}}ne choice, about the widest I \deleted{\unsure{c|w}}\unsure{w|c}ould think of, is: for every pair $(f,g)$ in $\underline{C}_n$ satisfying condition $(1)$ of (A), add one \deleted{arrow} new arrow $h$. To avoid set-theoretic difficulties though, we better \deleted{\unsure{reduce}} first modify the definition of $\underline{C}_{\unsure{o|0}}$ so that the set of its objects should be in the universe we are working in, preferably even it be countable. Or else and more reasonably, we will pic\handwritten{k} one $h$ for every isomorphism class of situations $(f,g)$ in $\underline{C}_n$. Another restriction to avoid too much redundancy -- this was the first definition \unsure{actu}ally that flipped to my mind the day before yesterday -- is to add a \uline{new} $h$ only when there is no ``old'' one, namely in $\underline{C}_n$, serving the same purpose. Then it came to my mind that there is a lot of redundancy still, thus there would be already in

% archive sheet number: 108
\scanpage{110}\mspage{\handwritten{$9'$}}

infinitely many ope\deleted{\unsure{a}}rations standing for the single ope\unsure{rat}ion $v\overset{i}{\circ}u$ say, which could be viewed in effect in terms of an arbitrary $n$-sequence \handwritten{$(n\geq2)$} of ``composable'' $i$-objects $u_1=u,u_{\handwritten{2}}=v,u_{\handwritten{3}},\ldots\,u_n$. The natural way to meet this ``objection'' would be to restrict to \deleted{a} pairs $(f,g)$ which cannot be factored non trivially through another object $K'$ as
\diagram{\begin{tikzcd}[ampersand replacement=\&,column sep=large]
K \arrow[r] \& K' \arrow[r,shift left=1.4ex,"{f',g'}"] \arrow[r,shift right=1.4ex] \arrow[r,phantom,"{\deleted{F_i}}" description] \& F_i\; .
\end{tikzcd}}
But even with such restrictions, there remain a lot of redundancies -- and this again seems to me in the nature of things, namely that there is no really natural, ``most economic'' way for achieving condition (A), by a stepwise construction meeting condition (B). For instance, in $\underline{C}_1$ already we will have \deleted{the} not merely the compositions $v\overset{i}{\circ}u$, but at the same time simultaneous compositions
\[
u_n\handwritten{\circ}u_{n-1}\handwritten{\circ}\ldots\handwritten{\circ}u_1 \tag*{(*)}
\]
for ``composable'' sequences of $i$-objects $(i\handwritten{\geq}1)$, without reducing this (as is c\unsure{o|u|a}stomary) to iteration of the binary composition $v\circ u$. Of course using the binary composition, and more generally iteration of $n'$-ary compositions with $n'\handwritten{<}n$ (when $n\handwritten{\geq}3$), we get an impressive bunch of ope\deleted{\unsure{a}}rations in the '$n$ variables $u_1,\ldots,u_n$, serving the same purpose as $(*)$. All these will be tied up by homotopies in the next step $\underline{C}_2$. We would like to think of this set of homotopies in $\underline{C}_2$ as a kind of ``transitive system of isomorphisms'' (of associativity), now the transitivity relations one is looking for will be replaced by homotopies again between compositions of homotopies, which will enter in the picture with $\underline{C}_{\handwritten{3}}$, \inserted{\handwritten{etc.}} Here the infinite story is exemplified by the more familiar situation of the two ways in which one could define a``$\unsure{\otimes|x}$-composition with associativity'' in a category, \deleted{using either} starting either in terms of a binary operation, or \inserted{with} a bunch of $n$-ary operations -- with\handwritten{,} in both ca\handwritten{\unsure{ses}}, the associativity isomorphisms being an essential part of the structure. Here again, while it is generally (and quite validly) felt that the two points of view are equivalent\unsure{;|,} and both have their advantages and their drawbacks, still it is not true \inserted{I believe} that the \inserted{two} categories of algebraic structures ``category with associative $\unsure{\otimes|x}$-operation'', using one or the other definition, are equivalent. Here the story though of the relation between the two notions is a finite one, due to the fact that it is related to the notion of $2$-categories or $2$-groupoids, instead of $\infty$-groupoids as before \ldots

Thus I dont feel really like spending much energy in cutting down redundanci\unsure{es}, but prefer working with a notion of $\infty$-groupoid which remain\unsure{s} partly indeterminate, the main features being embodied in the conditions (A) and (B) and in the description of $\underline{C}_{\unsure{o|0}}$, without other specification.

% archive sheet number: 109
\scanpage{111}\mspage{\handwritten{10}}

\handwritten{11.} One convenient way for constructing \deleted{\unsure{$\underline{C}$}} a category $\underline{C}_{\handwritten{\infty}}$ would be to defin\unsure{e} for every $K,L\handwritten{\in}\operatorname{Ob}(\underline{C}_{\unsure{o|0}})=\operatorname{Ob}(B_o)$ the set $\operatorname{Hom}_{\handwritten{\infty}}(K,L)$ as a subset of the set $\operatorname{Hom}(\handwritten{|}L\handwritten{|},\handwritten{|}K\handwritten{|})$ of continuous maps \deleted{from} between the geometric realisations of $L$ and $K$ in terms of gluing together cells $D_i$, the composition of arrows in $\underline{C}_{\handwritten{\infty}}$ being just composition of maps. This amounts to defining $\underline{C}_{\handwritten{\infty}}$ as \deleted{\unsure{a} subsidiary of} the dual of a category $\underline{B}_{\handwritten{\infty}}$ of topological description\deleted{s}. It will be sufficient to define for every cell $D_n$ and every $K$ a subset $\operatorname{Hom}_{\handwritten{\infty}}(D_n,\handwritten{|}K\handwritten{|})$ \inserted{of $\operatorname{Hom}(D_n,\handwritten{|}K\handwritten{|})$}, satisfying the two conditions:

a) stability \deleted{with}\inserted{by} compositions $D_n\handwritten{\to|}K\handwritten{|}\handwritten{\unsure{\to|-}} \handwritten{|}K'\handwritten{|}$, where $K\qquad K'$ is \deleted{a continuous}\inserted{an ``allowable''} map, \inserted{namely} subjected only to the condition that its restriction to any standart subcell $D_{n'}\handwritten{\unsure{\subset|<}|}K\handwritten{|}$ is again ``allowable'' i.e. in $\operatorname{Hom}_{\handwritten{\infty}}$.

b) Any ``allowable'' map $S_n^{\deleted{\unsure{+|\varepsilon}}}\handwritten{\to|}K\handwritten{|}$ (i.e. whose restrictions to $S_n^{\handwritten{+}}$ and $S_n^{\handwritten{-}}$ are allowable) extends to an allowable map $D_{n+1}\handwritten{\longrightarrow|}K\handwritten{|}$.

Condition a) merely ensures stability of allowable maps under composition, and the fact that $\underline{B}_{\handwritten{\infty}}$ (endowed with the allowable maps as morphisms) has the correct binary amalgamated sums, whereas b) expresses condition A) on $\underline{C}_{\handwritten{\infty}}$. These conditions are satisfied when we take as $\operatorname{Hom}_{\handwritten{\infty}}$ subsets defined by tameness conditions (such as piecewise linear for suitable piecewise linear structure\inserted{\handwritten{ in the $D_i$'s}}, or differentiable etc). The condition (B) however is of a subtler nature in this topological interpretation and surely not met by such sweeping tameness requirements only ! Finally the question as to whether we can actually in this way describe an ``acceptable'' category $\underline{C}_{\handwritten{\infty}}$, by defining sets $\operatorname{Hom}_{\handwritten{\infty}}$ namely describing $\underline{C}_{\handwritten{\infty}}$ in term\unsure{s} of $\underline{B}_{\handwritten{\infty}}$, seems rather subsidiary\inserted{\handwritten{ after all}}. We may think of course of constructing stepwise $\underline{B}_{\handwritten{\infty}}$ via subcategories $\underline{B}_n$, by adding stepwise new arrows in order to meet condition b), thus paraphrasing condition (B) for passage \deleted{of} from $\underline{C}_n$ to $\underline{C}_{n+1}$ -- but it is by no means clear that when passing to the \deleted{\unsure{sub}}category $\underline{B}_{n+1}$ by composing \deleted{arrows} maps of $\underline{B}_n$ and ``new'' ones, and using amalgamated sums too, there might not be some \inserted{undesirable} extra\deleted{\unsure{s}} relations \deleted{\unsure{arising}} in $\underline{B}_{n+1}$, coming from the topological interpretation of the arrows in $\underline{C}_{n+1}$ as maps. To say it differently, universal algebra furnishes us readily with an acceptable sequence of categories $\underline{C}_n$ and hence $\underline{C}_{\handwritten{\infty}}$, and by the universal properties of the $\underline{C}_n$ in terms of $\underline{C}_{n-1}$, we readily get (using arbitrary choices) a contravariant functor $K\handwritten{\to|}K\handwritten{|}$ from $\underline{C}_{\handwritten{\infty}}$ to the category of topological spaces (i.e. a co-$\infty$-groupoid in (Top)), but it is by no means clear that this functor is faithful -- and it does'nt really matter after all !

% archive sheet number: 110
\scanpage{112}\mspage{\handwritten{$10'$}}

\handwritten{12.} I think I really better stop now, except for one last comment. The construction of a co-$\infty$-groupoid in (Top), giving rise to the fundamental functor
\[
(\mathrm{Top})\handwritten{\longrightarrow}(\infty\text{-}\unsure{\text{grouppoids}|\text{groupoids}}),
\]
generalises,as I already alluded to earlier, to the case when (Top) is replaced by an arbitrary``model category''\inserted{ $M$} in your sense. Here however the choices occur not only \inserted{stepwise} for the primary, secondary, ternary etc structures, but already for the primitive structure, namely \inserted{\handwritten{by}} choice of objects $D_i$ $(i\handwritten{\in}\underline{N})$ in $M$, and source and target maps $D_{\unsure{1|i}\deleted{\unsure{-1}}}\handwritten{\rightrightarrows}D_{\unsure{o|0}}\deleted{\unsure{i\geq1}}$. These choices can be made inductively, by choosing first for $D_o$ the final object, or more generally any object which is fibrant and trivial (over the final object), $D_{-1}$ being the initial object, and defining further $S_o=D_o\handwritten{\amalg}_{D_{-1}}D_o=D_o\handwritten{\amalg}D_o$ with obvious maps ${\handwritten{\unsure{\kappa|\phi}}}_o^{+},\handwritten{\kappa}_o^{-}:D_o\handwritten{\rightrightarrows}S_o$, and then \deleted{if}if everything is constructed up to $D_{n\deleted{\unsure{1}}}$\deleted{\unsure{)}} and $S_n=(D_n,{\handwritten{\unsure{\phi|\delta}}}_{n-1}^{+})\mathop{\phantom{\amalg}}\limits_{D_{n-1}}(D_n,\handwritten{\phi}_{n-1}^{-})$, defining $D_{n+1}$ as any fibrant and trivial object together with a cofibrant map
\[
S_n\handwritten{\to}D_{n+1}\quad ,\qquad \unsure{\mathrm{e}|?}
\]
and $\handwritten{\phi}_n^+,\handwritten{\phi}_n^-$ as the compositions of the latter with $\handwritten{\kappa}_n^+,\handwritten{\kappa}_n^-:D_n\handwritten{\rightrightarrows}S_n$. Using this and amalgamated sums in $M$, we get our functor
\[
\handwritten{\underline{B}_o =} \underline{C}_o^{o}\handwritten{\to}M\quad ,\qquad K\handwritten{\mapsto|}K\handwritten{|}_M\quad ,
\]
commuting with amalgamated sums\unsure{,|.} which we can extend stepwise through the $\underline{C}_n^{o}$'s to a functor \inserted{\handwritten{$\underline{B}_{\infty}=$}}$\underline{C}_{\handwritten{\infty}}^{o}\handwritten{\to}M$, provided we know that the objects $\handwritten{|}K\handwritten{|}_M$ of $M$ $(K\handwritten{\in}\operatorname{Ob}\underline{C}_o)$ obtained \deleted{through} by ``standart'' gluing of the $D_n$'s in $M$, are again fibrant and trivial -- and I hope indeed that your axioms imply that, via,say, that if $Z\handwritten{\to}X$ and $Z\handwritten{\to}Y$ are cofi\handwritten{b}rant and $X,Y,Z$ are fibrant and trivial, then $X\handwritten{\amalg}_ZY$ is fibrant and trivial \ldots

Among the things to be checked is of course that when we localize the category of \deleted{\unsure{(}}$\infty$-groupoids with respect to morphisms which are ``weak equivalences'' in a rather obvious sense (NB the definition of the $\handwritten{\pi}_i$'s of an $\infty$-groupoid is practically trivial !), we get a category equivalent to the usual homotopy category (Ho). Thus we get a composed functor
\[
M\qquad (\infty\text{-groupoids})\qquad(\mathrm{Ho})\quad ,
\]
as announced. I have some intuitive feeling of what this functor stands for, at least when $M$ is say the category of semisimplicial sheaves, or (more or less equivalently) of $n$-\deleted{groupoids}\inserted{gerbs} or $\infty$-\deleted{groupoids}\inserted{gerbs} on a given topos: namely it \inserted{should} correspond\deleted{s} to the operation of ``integration'' or``secti\deleted{\unsure{o}}ons'' for $n$-gerbs (more generally for $n$-stacks) over a topos -- which is indeed \uline{the} basic operation (embodying non commutative cohomology objects of the topos) in ``non commutative homological algebra''.

I guess that's about it for today -- \unsure{I|i}t's getting late and time to go to bed ! Good night .

% archive sheet number: 111
\scanpage{113}\mspage{\handwritten{11}}

\hfill 22.2.19\unsure{3|8}\handwritten{3}\\

\handwritten{13.} It seems I can't help pursuing further the reflection I started with this letter ! First I would like to come back upon terminology. Maybe to give the name of $n$-groupoids and $\infty$-groupoids to the objects I was after is not proper, for two reasons: a) it conflicts with a standing terminology, applying to structure species which are frequently met and deserve names of their own, even if they turn out to be too restrictive kind of objects for the use I am having in mind -- so why not keep the terminology already in use, especially for two-groupoids, which is\deleted{not bad} pretty well sui\handwritten{t}ed after all; b) the structure species I have in mind is not really a very well determined one; it depends as we saw on choices, without any one choice looking convincingly better than the others -- so it would be a mess to give an unqualified name to such structure, depending on the choice of a certain category $\underline{C}_{\infty}=\underline{C}$. I have been thinking of the terminology $n$-\uline{stack} and $\infty$-\uline{stack} (stack = ``champ'' in French), a name introduced in Giraud's book (he was restricting to champs = $1$-champs), which over a topos reduced to a one-point space reduces in his case to the usual notion of a category i.e. $1$-category. Here of course we are thinking of ``stacks of groupoids'' rather than arbitrary stacks, which I would like to call (for arbitrary order $n\handwritten{\in}\underline{N}$ or $n=\infty$) $n$-Gr-stack\unsure{.|s} -- suggesting evident ties with the notion of Gr-categories, we should say Gr-1-categories, of Mme Hoang Xuan Sinh. One advantage of the \handwritten{\unsure{name}} stack\unsure{''|'} is that the use it had so far spontaneously suggests the extension of these notions to the corresponding notions over an arbitrary topos, which of course is what I am after ultimately. Of course, when an ambiguity is possible, we should speak of \deleted{\unsure{$\underline{C}$}} $n$-$\underline{C}$-stacks -- the reference to $\underline{C}$ should make superfluous the ``Gr'' specification. Thus $n$-$\underline{C}$-stacks are essentially the same as $n$-\deleted{\unsure{$\underline{C}$}}-stacks over the final topos i.e. over a one point space. When both $\underline{C}$ and ``Gr'' are understood in a given context, we will use the \deleted{word ``stack''} terminology $n$-stack simply, or even``stack'' when $n$ is fixed throughout. Thus it will occur that in certain contexts ``stack'' will just mean a usual groupoid, \deleted{or a 2-groupoid,} in others it will mean just a category, \deleted{or a 2-category} but when $n=2$ it will not mean a usual 2-groupoid, but something more general, defined in terms of $\underline{C}$.

The categories $\underline{C}=\underline{C}_{\infty}$ described before merit a name too -- I would like to call them ``\uline{coherators}'' (``\uline{cohéreurs}'' \unsure{in|en} French)\unsure{.|,} \deleted{in order to} \handwritten{T}his name is meant to suggest that $\underline{C}$ embodies a hierarchy of coherence relations, more accurately of coherence``homotopies''. When dealing with stacks, the term $i$-\uline{homotopies} (rather than $i$-objects or $i$-arrows) for the elements of the $i^{\mathrm{th}}$ component $F_i$ of a stack seems to me the most suggestive -- they will of course be denoted graphically by arrows,

% archive sheet number: 112
\scanpage{114}\mspage{\handwritten{$11'$}}

such as \unsure{(|h}$h:f\handwritten{\to}g$ in the formulation of (A)\unsure{,|'} yesterday.) More specifically, I will call coherator any category \uline{equivalent} to a category $\underline{C}_{\infty}$ as constructed before. \deleted{\unsure{Such}} Thus a coherator is stable under binary \deleted{amalgamated sums}\inserted{\handwritten{fiber products}}, moreover the $F_i$ are recovered up to isomorphism as the indecomposable elements of $\underline{C}$ with respect to amalgamation.
% slanted
\margin{\handwritten{false, see \unsure{PS|ps|P.S.}}}
However, in a category $\underline{C}_{\infty}$, the objects $\underline{F}_i$ have non trivial automorphisms -- namely the ``duality involutions'' and their compositions (the group of automorphisms of $\underline{F}_i$ should turn out to be canonically isomorphic to $(\pm1)^i$), in \deleted{\unsure{and the $F_i$}} other words by the mere category structure of a coherator we will not be able to recover the objects $\underline{F}_i$ in $\underline{C}$ up to unique isomorphism\unsure{.} Therefore in the structure of a coherator should be included too the choice of the basic undecomposable objects $\underline{F}_i$ (one in each isomorphism class), and moreover the arrows $\underline{s}_{\unsure{1|l}}^i,\underline{t}_1^i:\underline{F}_i\handwritten{\rightrightarrows}\underline{F}_{i-1}$ for $i\handwritten{\geq}1$ (a priori, only the \uline{pair} $(\underline{s}_1^i,\underline{t}_1^i)$ can be described intrinsically in terms of the category structure of $\underline{C}$, once $\underline{F}_i$ and $\underline{F}_{i-1}$ are chosen \ldots). But it now occurs to me that we dont have to put in this extra structure afte\unsure{r} all. -- while the $\underline{F}_i$'s separately do have automorphisms, the system of the objects $(\underline{F}_i)_{i\handwritten{\in}\underline{N}}$ and of the maps $(\underline{s}_1^i)_{i\handwritten{\geq}1}$ and $(\underline{t}_1^i)_{i\handwritten{\geq}1}$ has only the trivial automorphism (all this of course is heuristics, I did'nt really prove anything -- but the structure of the full subcategory of a $\underline{C}_{\infty}$ formed by the objects $\underline{F}_i$ seems pretty obvious \ldots). To finish getting convinced that the mere category structure of a coherator includes already all other relevant structure, we should still describe a suitable intrinsic filtration by subcategories $\underline{C}_n$. We define the $\underline{C}_n$ inductively, $\underline{C}_{\unsure{o|0}}$ being the``primit\deleted{\unsure{y}}i\inserted{ve} structure'' (the arrows are those deducible from the source and target arrows by composition and \deleted{\unsure{amalgamation}}\inserted{fiber products}), and $\underline{C}_{n+1}$ being defined in terms of $\underline{C}_{\unsure{n|1}}$ as follows: add to $\mathrm{Fl}\,\underline{C}_n$ all arrows of $\underline{C}$ of the type $h:K\handwritten{\to}\underline{F}_i$ $(i\handwritten{\geq}1)$ such that $\underline{s}_1^i h$ and $t_1^i h$ are in $\underline{C}_n$, and the arrows deduced from the bunch obtained \deleted{through} by composition and fibered products. In view of these constructions, it would be an easy exercise to give an \inserted{intrinsic} caracterisation of a coherator, as a category satisfying certain internal properties.

I was a little rash right now when making assertions about the structure of the group of auomorphisms of $\underline{F}_i$ -- I forgot that two days ago I pointed out to myself that even the basic operation $\mathrm{inv}_i$ upon $\underline{F}_i$ need not \inserted{\handwritten{even}} be \deleted{an} involution\handwritten{s}!
% slanted
\margin{\handwritten{not even automorphisms!}}
However I just checked that if in the inductive construction of coherators $\underline{C}_{\infty}$ given yesterday, we insist on the most trivial irredundancy condition \deleted{\unsure{--}}(namely that we dont add a``new'' homotopy $h:f\handwritten{\to}g$ when there is already an old one), then any morphism $h:\underline{F}_i\handwritten{\to}\underline{F}_i$ such that $\underline{s}\,f=\underline{s}$ and $\underline{t}\,f=\underline{t}$, is the identity -- and that

% archive sheet number: 113
\scanpage{115}\mspage{\handwritten{12}}

\inserted{\handwritten{implies}} inductively that an automorphism of the system of $\underline{F}_i$'s related by the source and target maps $\underline{s}_{\unsure{1|l}}^i,\underline{t}_1^i$ is the identity. Thus it is correct afte\unsure{r} all, it seems, that the category structure of a coherator implies all other structure relevant to us.
% slanted
\margin{\handwritten{Maybe false, \unsure{we} should give moreover the subcategory $\underline{C}_{\unsure{o|0}}$ perhaps}}

I do believe that the description given so far of what I mean by a coherator, \deleted{\unsure{acting a}} namely something acting like a kind of pattern in order to define a corresponding notion of ``stacks'' (which in turn should be the basic coefficient objects in non commutative homological algebra, as well as a convenient description of homotopy types) embodies some of the essential features of the theory\inserted{ still in embryo} that wants to be developped. It is quite possible of course that some features are lacking still, for instance that some extra conditions have to be imposed upon $\underline{C}$, possibly of a very different nature from mere irredundancy conditions (which, I feel, are kind of irrelevant in this set-up). Only by pushing ahead and working out at least in outline the main aspects of the formalism of stacks, will it become clear whether or not extra conditions on $\underline{C}$ are needed. I would like at least to make a commented list of these main aspects, and possibly do some heuristic pondering on some of these in the stride, or afterwards. For today it seems a little late though -- I have been pretty busy with non mathematical work most part of the day, and the next two days I'll be busy at the university. Thus I guess I'll send off this letter tomorrow, and send you \unsure{an}\inserted{ later} elaboration (presumably much in the style of this unending letter) if you are interested. In any case I would appreciate any comments you make -- that's why I have been writing you after all ! I will probably send copies to Ronnie Brown, Luc Illusie and Jean Giraud, in case they should be interested (I guess at least Ronnie Brown is). Maybe the theory is going to take off after all, in the long last !

Very cordially yours\\
\hfill\handwritten{Alexander}\\

PS (23.2) I noticed a rather silly mistake in the notes of two days ago,when stating that the categories $\underline{C}_n$ admit fiber\deleted{ed} products -- what is true is that the category $\underline{C}_o$ has fiber\deleted{ed} products (by construction\inserted{ practically}), and that these are fiber\deleted{ed} products also in the categories $\underline{C}_n$ (by construction equally) i.e. that the inclusion functors $\underline{C}_o\handwritten{\to}\underline{C}_n$ commute to fiber\deleted{ed} products. Stacks in a category $C$ correspond to functors $\underline{C}_{\infty}\handwritten{\to}C$ whose restriction to $\underline{C}_o$ commutes with fiber products. I carried the mistake along in the \deleted{notes} yesterday notes -- it does'nt really change anything substantially. I will have to come back anyhow upon the basic notion of \unsure{a coherator}

% archive sheet number: 114
\scanpage{128}
% printed form
% Exposed strip of the underlying listing, reading down the right edge.
\begin{center}
\begin{tabular}{r@{\qquad}l}
\unsure{J*} & \unsure{:}*\\
& *\\
& *\\
* & *\\
& *\\
& \unsure{*}\\
\end{tabular}
\end{center}

% Reverse-side impressions, read with the scan reflected top to bottom.
% The banner is clipped along its left edge.
\begin{center}\ttfamily
\begin{tabular}{l@{\qquad}l}
\unsure{TTTTTT} & \unsure{SSSSSSSSSS|5555555555}\\
\unsure{TTTTTT} & \unsure{SSSSSSSSSSSS}\\
\unsure{TT} & \unsure{SS\hspace{6em}SS}\\
\unsure{TT} & \unsure{SS}\\
\unsure{TT} & \unsure{SS}\\
\unsure{TT} & \unsure{SSSSSSSSS}\\
\unsure{TT} & \unsure{\hspace{1em}SSSSSSSSS}\\
\unsure{TT} & \unsure{\hspace{6em}SSS}\\
\unsure{TT} & \unsure{\hspace{7em}SS}\\
\unsure{TT} & \unsure{\hspace{7em}SS}\\
\unsure{TT} & \unsure{SSSSSSSSSSS}\\
\unsure{TT} & \unsure{SSSSSSSSSS}\\
\end{tabular}
\end{center}

\begin{center}\scriptsize\ttfamily
\begin{tabular}{@{}lllllll@{}}
\unsure{50 PM} & \unsure{29 OCT 82} & \unsure{B6.PR1|R6.PR1} & SYS JES2 & STC & \unsure{32} END & \unsure{P****}\\
\unsure{50 PM} & 29 OCT \unsure{82} & \unsure{B6.PR1} & SYS JES2 & STC & 32 END & \unsure{P****}\\
\unsure{50 PM} & \unsure{29 OCT 82} & \unsure{B6.PR1} & SYS JES2 & \unsure{STC} & \unsure{32} END & \unsure{P****}\\
\unsure{50 PM} & \unsure{29 OCT 82} & \unsure{B6.PR1} & \unsure{SYS JES2} & STC & \unsure{32} \unsure{END} & \unsure{P****}\\
\unsure{50 PM} & \unsure{29 OCT 82} & \unsure{B6.PR1} & SYS JES2 & STC & 32 END & \unsure{P****}\\
\unsure{50 PM} & \unsure{29 OCT 82} & \unsure{B6.PR1} & SYS JES2 & STC & 32 END & \unsure{P****}\\
\end{tabular}
\end{center}

% Lower edge of the exposed listing; the tops of the first group are covered.
\hfill\texttt{\unsure{JIMS|JMS}\qquad *//}

% archive sheet number: 128
% A second, very faint lower-left number may read 127.
